% problem-000073 7 钙钛矿衍生结构
\section*{7. 钙钛矿衍生结构}
\textit{下列为分问。}

\subsection*{7-1}
后续研究发现，Bartlett当时得到的并⾮ $\mathrm { X e ^ { + } [ P t F _ { 6 } ] ^ { - } }$ −，⽽可能是 $\mathrm { X e F ^ { + } [ P t _ { 2 } F _ { 1 1 } ] ^ { - } }$ 。

\subsection*{7-1}
-1 写出 $\mathrm { X e F ^ { + } [ P t _ { 2 } F _ { 1 1 } ] }$ −中Xe的氧化数。

\subsection*{7-1}
-2 在 $\mathrm { [ P t _ { 2 } F _ { 1 1 } ] }$ −结构中，沿轴向有四次轴。画出 $\mathrm { [ P t _ { 2 } F _ { 1 1 } ] }$ −的结构。

\subsection*{7-2}
后来，⼤量含 Xe—F 和 $\mathrm { X e - O }$ 键的化合物被合成出来，如 $\mathrm { X e O F _ { 2 \circ } }$ 。根据价层电⼦对互斥理论(VSEPR)，写出 $\mathrm { X e O F _ { 2 } }$ 的⼏何构型及中⼼原⼦所⽤杂化轨道类型。

\subsection*{7-3}
1974年合成了第⼀例含 $\mathrm { X e - N }$ 键的化合物： $\mathrm { X e F _ { 2 } }$ 和 $\mathrm { H N } ( \mathrm { S O _ { 2 } F } ) _ { 2 }$ 在 $0 ~ ^ { \circ } \mathrm { C }$ 的⼆氟⼆氯甲烷溶剂中按1:1的计量关系反应，放出HF得到⽩⾊固体产物A。A受热⾄ $7 0 ^ { \circ } \mathrm { C }$ 分解，A中的 $\mathrm { X e }$ ⼀半以 $\mathrm { X e }$ ⽓放出，其他两种产物与Xe具有相同的计量系数且其中⼀种是常⻅的氙的氟化物。写出A的化学式及其分解的反应⽅程式。

\subsection*{7-4}
1989年，随着超⾼压下Xe单质固体由⽴⽅密堆积变为六⽅密堆积结构并显⽰出⾦属性，Xe的“惰性”帽⼦被彻底摘掉。理论预测，⾼压下Xe可以参与形成更复杂的化合物，如 $\mathrm { C s ^ { I } X e ^ { I I } A u _ { 3 \circ } }$ 若 $\mathrm { C s ^ { I } X e ^ { I I } A u _ { 3 } }$ 采⽤钙钛矿类型的结构，分别写出 $\mathrm { C s ^ { I } }$ 和 $\mathrm { X e ^ { I I } }$ 最近邻的⾦原⼦数。
